\documentclass[11pt]{article}

\usepackage{amsmath}
\usepackage{amssymb}
\usepackage{amsthm}
\usepackage[margin=1in]{geometry}
\usepackage{hyperref}
\hypersetup{
    colorlinks=true,
    linkcolor=blue,
    citecolor=blue,
    urlcolor=blue
}

\title{Tutorial 7: Blockchain Mechanics and the Double-Spending Problem}
\author{Financial Economics \\ Research School of Economics, Australian National University \\
\small Instructor: Simon Mishricky}
\date{}

\begin{document}
\maketitle

\noindent\textbf{Source.} The problems below are based on Chiu and Koeppl (2017),
\emph{The Economics of Cryptocurrencies -- Bitcoin and Beyond} (cited throughout as CK2017).

\medskip

\noindent\textbf{Learning objectives.} After completing this tutorial you should be able to
\begin{itemize}
    \item explain how a blockchain functions as a decentralised record-keeping device and how
          Proof-of-Work (PoW) underpins consensus;
    \item articulate the double-spending problem and the role of confirmation lags;
    \item analyse the mining game and derive the no-double-spending constraint;
    \item compute the probability of a successful secret-mining attack as a function of trade size,
          rewards, and confirmation lag.
\end{itemize}

\bigskip

\begin{enumerate}

%--------------------------------------------------------------------
\item \textbf{Blockchain as memory.}
A central claim of CK2017 is that the blockchain is a record-keeping device.
\begin{enumerate}
    \item In one or two paragraphs, contrast a blockchain with traditional fiat money in terms
          of how each ``remembers'' past transactions. Why does this difference matter for
          decentralised payments?
    \item Why does the longest-chain rule make rewriting old transactions costly? Explain in terms
          of the PoW that must be redone if a malicious miner wants to revise a block $k$ periods
          in the past.
    \item Explain in your own words why a blockchain alone does \emph{not} eliminate
          forward-looking double-spending attacks.
\end{enumerate}

%--------------------------------------------------------------------
\item \textbf{The mining game.}
Following Section 3.2 of CK2017, $M$ symmetric miners each choose computing power $q(i)$ at unit cost.
The probability that miner $i$ wins block $n$ is
\[
    \rho(i) = \frac{q(i)}{\sum_{m=1}^M q(m)}.
\]
The winner receives reward $R$ (in real balances), discounted by $\beta$ across the period.
\begin{enumerate}
    \item Write down the miner's per-subperiod problem and derive the first-order condition for $q(i)$.
    \item Imposing symmetry $q(m)=Q$, show that the Nash equilibrium choice is
          \[
              Q \;=\; \frac{M-1}{M^2}\,\beta R.
          \]
    \item Compute the aggregate computing cost $MQ$ and show that as $M\to\infty$,
          $MQ \to \beta R$. Interpret this ``rent dissipation'' result.
    \item Compute the expected profit of each miner across the full period of $\bar N+1$
          subperiods and show that it vanishes in the competitive limit.
\end{enumerate}

%--------------------------------------------------------------------
\item \textbf{Double-spending attack mechanics.}
Consider the secret-mining attack of Section 3.3 of CK2017. The buyer must mine $N+1$ blocks
faster than the rest of the network. Let $\Delta = d/R + (N+1)$.
\begin{enumerate}
    \item Explain the intuition for why the buyer must mine $N+1$ \emph{consecutive} blocks
          (not $N$).
    \item In the final round (subperiod $N$), the buyer's payoff from investing $q_N$ is
          \[
              \rho(q_N)\,\beta[d + (N+1)R] - q_N,
          \]
          where the second term in brackets reflects the $N+1$ block rewards the buyer would
          claim after broadcasting his secret chain. Take $M\to\infty$. Show that the optimal
          investment satisfies
          \[
              \hat q_N(d,N) = \beta R\bigl[\sqrt\Delta - 1\bigr]
          \]
          and that the corresponding success probability is
          \[
              \rho(\hat q_N) = \frac{\sqrt\Delta - 1}{\sqrt\Delta}.
          \]
    \item \emph{Backward induction.} Using the recursive structure
          $D_n(d,N) = \max_{q_n}\rho(q_n)D_{n+1}(d,N) - q_n$ with terminal condition
          $D_N(d,N) = \beta R(\sqrt\Delta - 1)^2$, verify by induction that
          \[
              D_{N-s}(d,N) = \beta R[\sqrt\Delta - (s+1)]^2,\quad s=0,1,\ldots,N.
          \]
\end{enumerate}

%--------------------------------------------------------------------
\item \textbf{The no-DS constraint.}
\begin{enumerate}
    \item Use the result $D_0(d,N) = \beta R[\sqrt\Delta - (N+1)]^2$ to derive the
          ``double-spending proof'' condition $d < R(N+1)N$.
    \item Show that the unconditional probability of a successful double spend (when the
          DS constraint fails) is
          \[
              P(d,N) = \frac{\sqrt\Delta - (N+1)}{\sqrt\Delta}.
          \]
    \item Discuss qualitatively how $P(d,N)$ varies with $d$, $R$, and $N$. In particular,
          explain why doubling the confirmation lag has a non-linear effect on attack probability.
    \item Suppose $R = 0.5$ and consider $N=6$ (Bitcoin's typical merchant standard).
          What is the largest payment $d$ that can be settled with full finality? Now repeat
          for $N=1$. Comment on the trade-off.
\end{enumerate}

%--------------------------------------------------------------------
\item \textbf{Immediate vs.\ final settlement.}
CK2017 prove a sharp impossibility theorem: under any PoW protocol, settlement cannot be
both immediate and final.
\begin{enumerate}
    \item Show, using your derivations above, that at $N=0$ the DS constraint
          $d < R(N+1)N$ fails for any $d>0$.
    \item Explain in plain English what this implies for retail spot trade in a Bitcoin-like
          system.
    \item Compare and contrast the cryptocurrency trade-off (security vs.\ speed of settlement)
          with the trade-off faced by a centralised payment system (e.g.\ Visa, Fedwire).
\end{enumerate}

%--------------------------------------------------------------------
\item \textbf{Mining-reward composition (short).}
The mining reward in the general equilibrium model is
\[
    R = \frac{Z(\mu - 1) + D\tau}{\bar N+1},
\]
where $\mu$ is the gross money growth rate, $\tau$ the transaction fee, $Z$ aggregate balances,
and $D$ aggregate transfers.
\begin{enumerate}
    \item Explain the economic interpretation of each of the two terms in the numerator.
    \item Why is $R$ divided by $\bar N+1$?
    \item Discuss why, holding $R$ fixed, the choice between financing rewards via $\mu$ or $\tau$
          might affect different buyers asymmetrically. (We will return to this in Tutorial 8.)
\end{enumerate}

\end{enumerate}

\begin{thebibliography}{9}
\bibitem{chiu2017}
Chiu, J.\ and T.\ V.\ Koeppl (2017). The Economics of Cryptocurrencies -- Bitcoin and Beyond.
\textit{Bank of Canada Staff Working Paper / Queen's University}.

\bibitem{lagos2005}
Lagos, R.\ and R.\ Wright (2005). A Unified Framework for Monetary Theory and Policy Analysis.
\textit{Journal of Political Economy} 113(3), 463--484.
\end{thebibliography}

\end{document}
